\subsection{Square-free polynomials as a finite-root interface}

Let $\mathcal P_n\cong\C^n$ be the space of monic degree-$n$ complex
polynomials and let $\Delta_n\subset\mathcal P_n$ be its algebraic discriminant.
The square-free locus is
\[
 \sfP_n=\mathcal P_n\setminus\Delta_n.
\]
Vieta's map identifies it with the unordered configuration space
\begin{equation}
 \sfP_n\cong\UConf_n(\C)
 =\Conf_n(\C)/S_n.
 \label{eq:vieta-configuration}
\end{equation}
Its fundamental group is the Artin braid group $B_n$
\cite{Artin1925TheorieDerZoepfe,FadellNeuwirth1962ConfigurationSpaces,Birman1974Braids}.

\begin{theorem}[Finite-root incidence and braid monodromy]
\label{thm:finite-root-braid}
Let $B$ be a connected smooth manifold and let
\begin{equation}
 F_b(w)=w^n+c_1(b)w^{n-1}+\cdots+c_n(b)
 \label{eq:polynomial-family}
\end{equation}
be a smooth family of monic polynomials with nonvanishing discriminant.  Then
\begin{equation}
 \calR_F=\{(w,b)\in\C\times B:F_b(w)=0\}
 \label{eq:root-incidence}
\end{equation}
is a smooth real codimension-two submanifold whose
projection to $B$ is a proper $n$-sheeted covering.  The coefficient map
$c:B\to\sfP_n$ induces
$c_*:\pi_1(B,b_0)\to\pi_1(\sfP_n,c(b_0))$.  After choosing the usual path to a
standard base configuration and the resulting identification of the target with
$B_n$, this is braid monodromy
\begin{equation}
 c_*:\pi_1(B,b_0)\longrightarrow B_n.
 \label{eq:braid-monodromy-map}
\end{equation}
Changing that identification conjugates the representation.
Its composition with $B_n\to S_n$ is the ordinary permutation monodromy of the
covering $\calR_F\to B$.
\end{theorem}

\begin{proof}
At a root of a square-free polynomial,
$\partial F_b/\partial w\ne0$.  The implicit-function theorem therefore writes the
root incidence locally as $n$ disjoint graphs over $B$, so it is an $n$-sheeted
covering.  A finite covering of locally compact Hausdorff spaces is proper.  The
map to unordered root configurations is exactly the coefficient map under
\eqref{eq:vieta-configuration}; functoriality of the fundamental group gives
\eqref{eq:braid-monodromy-map}.  Forgetting paths and retaining only their endpoint
labels is the standard quotient $B_n\to S_n$.
\end{proof}

For $B=S^1$, the incidence is a closed geometric braid in
$\C\times S^1$.  A basepoint and an ordering above it give an element of $B_n$;
changing the basepoint or ordering conjugates that element.  This is strictly more
information than the endpoint permutation.

\begin{warning}[Rank-one assignments do not yield $B_n$]
\label{warn:rank-one-not-braid}
The field $F$ in Theorem~\ref{thm:finite-root-braid} is complex-valued and has real
rank two.  The theorem does not turn a real AEG assignment into a polynomial, nor
does it promote a moving family of zero curves to a finite-root braid.  The two
constructions share the parameterized zero-section theorem but have different
codimensions.
\end{warning}

\subsection{Arithmetic-holomorphic braid realization}

On the complete basic hyperbolic AES with $\mu,\lambda>0$, Paper II constructs a global coordinate
$w=X+iy\in\HH$ in which arithmetic holomorphicity is ordinary holomorphicity.
This supplies a precise AEG background for Theorem~\ref{thm:finite-root-braid}.

\begin{theorem}[AEG realization of every braid]
\label{thm:aeg-braid-realization}
For every $\beta\in B_n$ and every complete basic hyperbolic AES with
$\mu,\lambda>0$, there is a smooth loop of arithmetic-holomorphic fields
$F_\theta$ such that:
\begin{enumerate}
\item every $F_\theta$ is a monic degree-$n$ polynomial in the arithmetic
holomorphic coordinate $w$;
\item all roots lie in a fixed disc $D\Subset\HH$ and remain simple;
\item the root incidence realizes the geometric braid $\beta$.
\end{enumerate}
Consequently, every oriented link is the closure of the root incidence of such an
arithmetic-holomorphic family.
\end{theorem}

\begin{proof}
Represent $\beta$ by a smooth loop
$\{z_1(\theta),\ldots,z_n(\theta)\}$ in
$\UConf_n(\C)$.  On the parameter interval $[0,2\pi]$, choose a smooth ordered
lift $(z_1(\theta),\ldots,z_n(\theta))$; its two endpoints may differ by the
permutation underlying $\beta$.  An affine translation and positive rescaling
place the compact unordered image in a disc $D\Subset\HH$ without changing its
braid class.  Set
\begin{equation}
 F_\theta(w)=\prod_{j=1}^n\bigl(w-z_j(\theta)\bigr).
 \label{eq:aeg-braid-polynomial}
\end{equation}
Because the endpoint lists differ only by a permutation, the coefficients are
smooth and periodic in $\theta$.  The roots are simple, and
each slice is holomorphic in $w$.  Paper II identifies it as an
arithmetic-holomorphic field.  Its root incidence is the chosen loop of
configurations.  The last statement follows from Alexander's closed-braid theorem
\cite{Alexander1923KnottedCurves}.
\end{proof}

The theorem is a universality result, not an invariant.  Because every braid can
be realized, the predicate ``admits an AEG root-tube realization'' cannot
distinguish links.  New information would have to be supplied by a natural
history-derived decoration or by a restricted class of allowed families.

\subsection{Logarithmic lifting of root tubes}

Suppose all roots avoid the origin.  Pull the root incidence back along
$\exp:\C\to\C^\times$:
\begin{equation}
 \widetilde{\calR}_F
 =\{(b,\omega):F_b(e^\omega)=0\}.
 \label{eq:log-root-incidence}
\end{equation}

\begin{theorem}[Logarithmic root transport and gauge]
\label{thm:log-root-transport}
The map $\widetilde{\calR}_F\to\calR_F$ is the canonical principal
$2\pi i\Z$-cover obtained by pullback of the exponential cover.  Its specified
deck-translation action is $\omega\mapsto\omega+2\pi i k$.  Along a parameter
loop, choose a labeling of the initial roots and
logarithmic lifts.  If root continuation has permutation $\sigma$, then there are
integers $k_i$ such that
\begin{equation}
 \omega_i(1)=\omega_{\sigma(i)}(0)+2\pi i k_i.
 \label{eq:log-endpoint-shift}
\end{equation}
Changing the initial lifts by
$\omega_i(0)\mapsto\omega_i(0)+2\pi i m_i$ changes the integers by
\begin{equation}
 k_i\longmapsto k_i+m_i-m_{\sigma(i)}.
 \label{eq:log-gauge-law}
\end{equation}
For every cycle of $\sigma$, the sum of the $k_i$ on that cycle is gauge invariant.
Relabeling the roots permutes these cycle sums, so their multiset, rather than an
ordered list, is label-independent.  The total sum $\sum_i k_i$ is gauge and label
invariant and equals the winding of the product of the roots around zero.
\end{theorem}

\begin{proof}
Equation~\eqref{eq:log-root-incidence} is the pullback of the exponential cover,
so the first claim is immediate.  Unique path lifting gives
\eqref{eq:log-endpoint-shift}.  If the initial lift of the $i$th path is shifted by
$2\pi im_i$, its entire lift is shifted by that amount.  Comparing the new endpoint
with the new lift assigned to the root $\sigma(i)$ gives
\eqref{eq:log-gauge-law}.  The $m_i$ terms telescope around each permutation cycle.
Finally, the product of all roots is, up to sign, the constant coefficient of
$F_b$; summing logarithmic increments computes its winding.
\end{proof}

Thus a vector $(k_i)$ is not an invariant before a lift gauge is fixed.  Cycle sums
are invariant but abelian.  The full noncommutative information remains in the
configuration path, more precisely in the annular braid group when roots are
constrained to $\C^\times$.

\subsection{Boundary braid and discriminant filling}

Theorem~\ref{thm:finite-root-braid} applies only in the complement of
$\Delta_n$.  A filling disc may meet the discriminant.  At a transverse simple
intersection, Proposition~\ref{prop:simple-discriminant-normal-form} gives
$w^2=\tau$, and its boundary loop contributes a conjugate of a standard positive
or negative half-twist according to orientation.  The collision is therefore a
two-cell event in the filling, while the boundary braid stays collision-free.

This separation resolves an otherwise persistent ambiguity:
\begin{center}
\small
\begin{tabular}{lll}
\toprule
Object & Allowed local event & Output \\
\midrule
regular real assignment family & no critical zero & curve tube \\
singular real family & Morse critical zero & birth or reconnection \\
square-free boundary root family & no repeated root & braid monodromy \\
root filling through $\Delta_n$ & simple double root & branch half-twist \\
\bottomrule
\end{tabular}
\end{center}

\begin{problem}[Natural braid monodromy]
\label{prob:natural-braid-monodromy}
Find a functorial procedure that assigns a finite complex-root field to an
expression history or to a real zero tube, and prove that the resulting braid is
independent of auxiliary coordinates and probes.  The arbitrary polynomial
realization of Theorem~\ref{thm:aeg-braid-realization} does not solve this problem.
\end{problem}
